% CSE558 Midterm Project Report -- IPL Impact Player study
% Compile with XeLaTeX (needs the Carlito font, as in the course template):  xelatex ipl_report.tex
\documentclass[10pt,a4paper]{article}
\usepackage[a4paper,top=1.6cm,bottom=1.8cm,left=1.7cm,right=1.7cm]{geometry}
\usepackage{fontspec}
\setmainfont{Carlito}
\usepackage{amsmath,amssymb}
\usepackage{graphicx}
\usepackage{booktabs,tabularx,array}
\usepackage{multicol}
\usepackage{enumitem}
\usepackage{titlesec}
\usepackage{caption}
\usepackage{xcolor}
\usepackage{needspace}
\usepackage[hidelinks]{hyperref}

\setlength{\columnsep}{0.9cm}
\setlength{\parindent}{0pt}
\setlength{\parskip}{4pt}
\linespread{1.03}
\titleformat{\section}{\Large\bfseries}{\thesection.}{0.5em}{}
\titlespacing*{\section}{0pt}{12pt}{5pt}
\titleformat{\subsection}{\normalsize\bfseries}{}{0em}{}
\titlespacing*{\subsection}{0pt}{8pt}{3pt}
\titleformat{\subsubsection}{\normalsize\bfseries}{}{0em}{}
\titlespacing*{\subsubsection}{0pt}{5pt}{2pt}
\captionsetup{font=small,labelfont=bf,skip=3pt}
\setlist{nosep,leftmargin=1.2em}
\newcolumntype{L}{>{\raggedright\arraybackslash}X}
\newcolumntype{P}[1]{>{\raggedright\arraybackslash}p{#1}}
\renewcommand{\arraystretch}{1.12}
\pagestyle{plain}

% compact hypothesis block: H0, H1, alpha, side, test, smallest difference
\newcommand{\hyp}[6]{%
\par\smallskip{\small
\begin{tabularx}{\columnwidth}{@{}>{\bfseries}p{0.40\columnwidth}L@{}}
\toprule
$\boldsymbol{H_0}$ & #1\\
$\boldsymbol{H_1}$ & #2\\
Significance level & #3\\
Side & #4\\
Test & #5\\
Smallest difference that matters & #6\\
\bottomrule
\end{tabularx}}\par}
% compact result block
\newcommand{\res}[6]{%
\par\smallskip{\small
\begin{tabularx}{\columnwidth}{@{}>{\bfseries}p{0.40\columnwidth}L@{}}
\toprule
Estimate & #1\\
95\% confidence interval & #2\\
Test statistic (d.f.) & #3\\
p-value (raw / Holm) & #4\\
Effect size & #5\\
Decision & #6\\
\bottomrule
\end{tabularx}}\par}
\newcommand{\figw}[3]{\par\smallskip\noindent\begin{minipage}{\columnwidth}\centering\includegraphics[width=#1\columnwidth]{#2}\captionof{figure}{#3}\end{minipage}\par\smallskip}
\newcommand{\fig}[2]{\par\smallskip\noindent\begin{minipage}{\columnwidth}\includegraphics[width=\columnwidth]{#1}\captionof{figure}{#2}\end{minipage}\par\smallskip}

\begin{document}

{\small\bfseries CSE558 Data Science \,$\bullet$\, Midterm Project Report \,$\bullet$\, Monsoon 2026}\par\vspace{6pt}
{\LARGE\bfseries Did the Impact Player rule change the IPL? Evidence from 19 seasons of ball-by-ball data}\par
\vspace{8pt}\hrule\vspace{6pt}

{\small
\begin{tabularx}{\textwidth}{@{}>{\bfseries}p{2.3cm}L@{}}
Group name & Null and Void\\
Dataset & \href{https://www.kaggle.com/datasets/abirambharathisa/ipl-ball-by-ball-2008-2026-all-19-seasons}{kaggle.com/datasets/abirambharathisa/ipl-ball-by-ball-2008-2026-all-19-seasons}\\
Code & \href{https://github.com/Snehil-Modi04/DSC_Project_IPL_Dataset}{github.com/Snehil-Modi04/DSC\_Project\_IPL\_Dataset} (notebook \texttt{ipl\_hypothesis\_tests.ipynb} and README)\\
\end{tabularx}

\vspace{4pt}
\begin{tabularx}{\textwidth}{@{}p{0.3cm}p{2.6cm}p{1.5cm}p{4.0cm}L@{}}
\toprule
\textbf{\#} & \textbf{Name} & \textbf{Roll no.} & \textbf{Institute email} & \textbf{Contribution}\\
\midrule
1 & Snehil Modi & 2024558 & snehil24558@iiitd.ac.in & Dataset collection, EDA, plot analysis, problem statements, report, code\\
2 & Vasu Mehta & 2024610 & vasu24610@iiitd.ac.in & Plot analysis, problem statements, H1, H2, report, code\\
3 & Ishaan Dhull & 2024267 & ishaan24267@iiitd.ac.in & Problem statements, H3, code\\
4 & Harshita Garg & 2024251 & harshita24251@iiitd.ac.in & H4, robustness summary, code\\
5 & Ishika Dhama & 2024270 & ishika24270@iiitd.ac.in & H5, limitations, code\\
\bottomrule
\end{tabularx}}

\vspace{4pt}
\begin{multicols}{2}
\raggedcolumns

\section{Dataset}
\textbf{Source.} Kaggle, ``IPL Ball-by-Ball 2008--2026: All 19 Seasons'', built from Cricsheet (\href{https://cricsheet.org}{cricsheet.org}) scorecards; downloaded 26 October 2026. Two linked tables: \texttt{ipl\_ball\_by\_ball\_2008\_2026.csv} (one row is one delivery) and \texttt{ipl\_matches\_2008\_2026.csv} (one row is one match: toss, result, method, totals), joined on \texttt{match\_id}. It covers all 1,243 IPL matches from 18 April 2008 to the 2026 final. We checked it against two other public IPL releases: every innings total agrees on all 1,095 and 1,169 shared matches. A fourth Kaggle ``IPL 2008--2025'' file we screened is synthetic (invented player names, wrong champions, $\sim$116 balls per match) and was rejected.

\textbf{What is raw about it.} 11,075 deliveries share a (match, innings, over, ball) key; 175 super-over balls sit beside regulation balls; wicket and extras fields are empty on 94--96\% of rows; DRS-review columns are 99.7\% empty; 60 matches were washed out, tied or shortened.

{\small
\begin{tabularx}{\columnwidth}{@{}Lrr@{}}
\toprule
 & \textbf{Raw file} & \textbf{After encoding}\\
\midrule
Instances (rows) & 295,732 & 295,557\\
Attributes (columns) & 38 & 114\\
Instances $\times$ attributes & 11,237,816 & 33,693,498\\
\bottomrule
\end{tabularx}}

\textbf{Size check.} $295{,}557 \times 114 = 33{,}693{,}498 \geq 3{,}000{,}000$: met with a margin of $\times 11.2$, and the raw ball table alone already gives 11.2M. One-hot encoding of batting and bowling team (15 each), venue (37), phase, extra type and dismissal kind takes the attributes from 38 to 114.

\subsection{Key features}
{\small
\begin{tabularx}{\columnwidth}{@{}P{2.45cm}P{1.5cm}L@{}}
\toprule
\textbf{Feature} & \textbf{Type} & \textbf{Real-world meaning}\\
\midrule
\texttt{total\_runs} & discrete & runs off one delivery (bat + extras)\\
\texttt{is\_six}, \texttt{is\_wicket} & binary & six hit / batter dismissed on this ball\\
\texttt{phase} & nominal & Powerplay (overs 1--6), Middle (7--15), Death (16--20)\\
\texttt{innings} & nominal & 1 = setting a target, 2 = chasing\\
\texttt{season} & discrete & 2008--2026; 2023+ is the Impact era\\
\texttt{toss\_winner}, \texttt{toss\_decision} & nominal & who won the toss; chose to bat or field\\
\texttt{winner} & nominal & match winner; outcome for H3--H5\\
first-innings total & discrete & derived from the balls; response for H1\\
\bottomrule
\end{tabularx}}

\section{Data preparation}
{\small
\begin{tabularx}{\columnwidth}{@{}P{2.3cm}P{2.3cm}L@{}}
\toprule
\textbf{Issue found} & \textbf{What we did} & \textbf{Why, and not the alternative}\\
\midrule
11,075 repeated ball keys & kept; new key \texttt{delivery\_\allowbreak in\_\allowbreak over} (0 repeats) & they are wides and no-balls, which are re-bowled with the same ball number; dropping them as duplicates would delete 11,075 real deliveries and their runs\\
175 super-over balls & removed & a tie-breaker, not part of either side's 20 overs; keeping them inflates totals\\
\bottomrule
\end{tabularx}}

{\small
\begin{tabularx}{\columnwidth}{@{}P{2.3cm}P{2.3cm}L@{}}
\toprule
\textbf{Issue found} & \textbf{What we did} & \textbf{Why, and not the alternative}\\
\midrule
extras and wicket fields empty (279,505) & filled with \texttt{none} & empty means ``no extra / no wicket'', not unknown (Fig.~\ref{fig:miss}); imputing a value would invent events\\
chase fields empty (52\%) & left empty, flagged as innings 1 & a first innings has no target; filling one would invent it\\
DRS columns 99.7\% empty & dropped & DRS began in 2018 and is used on 990 balls in total\\
\bottomrule
\end{tabularx}}

{\small
\begin{tabularx}{\columnwidth}{@{}P{2.3cm}P{2.3cm}L@{}}
\toprule
\textbf{Issue found} & \textbf{What we did} & \textbf{Why, and not the alternative}\\
\midrule
165 run rates empty & set to 0 & the over began with a wide: 0 legal balls, rate undefined\\
8 balls with 6 runs, not a six & used the \texttt{is\_six} flag & all-run 6s and overthrows are not sixes\\
19 ``retired hurt'' as wickets & not counted & the batter was not dismissed\\
franchise renames & verified one name per franchise (15) & otherwise Delhi Daredevils and Delhi Capitals split one history\\
60 matches: 9 no result, 16 ties, 23 D/L, 11 reduced, 2 curtailed & left out of score and chase tests & a shortened game is not a like-for-like 20-over contest\\
25 totals beyond the IQR fences & kept & the record 287 is real and is the effect H1 looks for; capping would erase it\\
\bottomrule
\end{tabularx}}

\fig{fig0_missing.pdf}{\label{fig:miss}Empty cells are structural, not lost data: almost every ball has no extra and no wicket, and only second innings have a target.}

\textbf{Unit of analysis.} Balls in one innings share batters, bowlers, pitch and match state, so they are not independent observations. No test runs on balls: H1, H2 and H5 use one first innings per match; H3 and H4 use one match.

\section{What changed after preparation}
{\small
\begin{tabularx}{\columnwidth}{@{}Lrr@{}}
\toprule
\textbf{Measure} & \textbf{Before} & \textbf{After}\\
\midrule
Ball rows & 295,732 & 295,557\\
Attributes & 38 & 114\\
Repeated ball keys & 11,075 & 0\\
Super-over balls & 175 & 0\\
Empty extras / wicket cells & 279,505 & 0\\
Innings totals disagreeing with match file & 0 & 0\\
Matches usable for score tests & 1,243 & 1,183\\
\bottomrule
\end{tabularx}}

\textbf{What we gained.} A unique key, zero ambiguous empty cells, and innings totals rebuilt from the balls that match the match file exactly. No value was invented: every step was a recode, a documented flag or a removal. \textbf{What we lost.} Only 175 super-over balls (0.06\%), plus 60 shortened matches (4.8\%) from the score tests.

\fig{fig_box_prep.pdf}{\label{fig:box}(a) First-innings totals before and after removing shortened matches: the lowest totals (down to 25) were rain-affected games. (b) Clean totals by era; dotted lines are the pooled IQR fences, numbers count totals beyond them.}
\textbf{What the box plots show.} (a) The 60 shortened matches have a median first-innings total of 150 against 168 for full-length games, and they hold the lowest score in the file (25). Leaving them in would drag the 20-over comparisons down with games that were never 20 overs long, so they leave the score and chase tests. (b) After that, 25 totals still lie beyond the IQR fences (86.8 and 252.8 runs). We keep them: they are real scores (the 287 is a record), and 13 of the 14 high ones come from 2023--26. Capping them would remove exactly the change H1 tests. The Impact-era box sits higher with a similar spread, which is why H1 compares means with Welch's t-test.

\section{Basic questions and problem statement}
\fig{fig1_seasons.pdf}{Season means with 95\% bands. Totals rose from 163 (2008) to 196 (2026), with a visible step after the rule; sixes follow the same path.}
\fig{fig2_totals.pdf}{First-innings totals, 2018--22 vs 2023--26: the whole distribution shifts right (means 168 and 190).}

\subsection{Basic questions}
\begin{enumerate}
\item Which franchise has won the most matches? (Mumbai Indians, 157.)
\item Do captains prefer to field after winning the toss? (66\% overall, 77\% since 2023.)
\item Are totals rising over the seasons? (Yes: Spearman $\rho = 0.84$ between season and mean total.)
\end{enumerate}

\subsection{Our problem statements}
\noindent\fbox{\parbox{\dimexpr\columnwidth-2\fboxsep-2\fboxrule}{\small
\textbf{P1 (primary).} Did the Impact Player rule (2023 on), which lets each side bring in a 12th player mid-match, raise scoring beyond the trend that already existed? (H1 totals, H2 sixes)\\[2pt]
\textbf{P2.} Do the toss and the choice to chase decide matches, and has the chasing edge moved now that targets are higher? (H3, H4)\\[2pt]
\textbf{P3.} Does losing 3+ wickets in the powerplay sink the batting-first side, and does deeper batting soften that? (H5)}}

\textbf{How they are novel.} The basic questions count wins and totals. Ours treat a rule change as a natural experiment, separate it from a pre-existing trend, and turn ball-level events (sixes, powerplay wickets) into innings-level outcomes that no single row records.

\textbf{Association or cause?} The rule was not randomised, so an era comparison alone cannot isolate it: bats, venues and the 2024 two-bouncer rule changed too. We therefore (i) model the trend and test the jump beyond it, (ii) compare the same venues before and after, (iii) run placebo ``rules'' in years with no rule. A jump that survives all three is strong evidence, but still not proof, of a rule effect.

\section{Statistical inferences}
\textbf{Rows tested.} Eras: \emph{Pre} = 2018--22 (303 clean matches), \emph{Impact} = 2023--26 (279), \emph{Early} = 2008--17 (601). H1, H2, H4, H5 use the 1,183 full-length matches; H3 uses all 1,218 matches with a winner.

\textbf{Design safeguards.} The five hypotheses, their sides and their real-world thresholds were written down before any outcome in the data was examined, so no hypothesis was chosen to fit a pattern we had seen. Each test is then challenged by the strongest alternative explanation we could think of (trend, venue, innings length, self-selection), and all five p-values carry a Holm correction. Every test reports an effect size, and every null result a power or equivalence analysis (summary in Table~\ref{tab:sum}).

\textbf{Why these tests.} Every primary test is a standard test from the course, chosen by the type of response and the number of groups. The regressions and the equivalence test appear only as robustness checks.
{\small
\begin{tabularx}{\columnwidth}{@{}P{0.5cm}P{1.75cm}P{1.65cm}L@{}}
\toprule
\textbf{H} & \textbf{Data} & \textbf{Test} & \textbf{Why this one}\\
\midrule
H1 & continuous; 2 independent groups & Welch two-sample t & $\sigma$ unknown and spreads unequal (Levene); with 2 groups it is one-way ANOVA ($F = t^2$)\\
H2 & counts, mildly skewed; 2 groups & Welch two-sample t & mean of $\sim$290 innings per group is near-normal (CLT); spreads unequal (Levene)\\
H3 & binary; 1 group vs 50\% & one-proportion z & one proportion against a fixed value, $np = 609 \ge 10$ so the normal approximation holds; same as $\chi^2$ goodness of fit\\
H4 & binary; 3 groups & $\chi^2$ independence & more than 2 groups, so a z-test cannot compare them at once\\
H5 & binary; 2 groups & two-proportion z & one-sided $H_1$; the $2\times2$ $\chi^2$ is only two-sided\\
\bottomrule
\end{tabularx}}

\textbf{How the chi-square and the other tests link.} Several tests are the same evidence in different forms, and the numbers confirm it. H1: one-way ANOVA gives $F(1, 580) = 67.0$, the square of the pooled $t = 8.18$. H3: $z^2 = 1.09^2 = 1.19$, exactly the $\chi^2$ goodness-of-fit statistic. H5: $z^2 = (-6.31)^2 = 39.8$, exactly the $2\times2$ $\chi^2$, and Fisher's exact test agrees. H4: the likelihood-ratio $G = 0.44$ matches $\chi^2 = 0.44$. So where we prefer z or t over $\chi^2$, it is for the one-sided alternative and the signed estimate, not because the tests disagree.

\subsection{H1 (P1): are first-innings totals higher in the Impact era?}
\subsubsection{a) Hypotheses, written before running the test}
\hyp{$\mu_I - \mu_P \le 0$}{$\mu_I - \mu_P > 0$}{$\alpha = 0.05$}{one-sided (right)}{Welch t-test}{10 runs}
$\mu_I, \mu_P$: mean first-innings total, Impact and Pre eras. One-sided because P1 claims a rise. 10 runs is roughly one extra boundary-over; below that, the par score a captain plans for does not change.

\subsubsection{b) Other tests we considered}
Two independent groups, a continuous response and an unknown $\sigma$ call for a two-sample t-test. We use Welch's version because the spreads differ: the SDs are 35.7 and 29.9, and Levene's test rejects equal variances ($p = 0.002$).
{\small
\begin{tabularx}{\columnwidth}{@{}P{2.3cm}L@{}}
\toprule
\textbf{Also considered} & \textbf{Why we did or did not use it}\\
\midrule
Pooled t-test & assumes equal variances, rejected above; gives $t = 8.18$ (580), same decision\\
Mann--Whitney / permutation & no normality needed; same decision (HL shift $+22$, 17 to 28; permutation $p < 10^{-4}$)\\
z-test & needs the population $\sigma$\\
OLS with trend (+ venue FE) & answers a sharper question, ``beyond the trend?''; run as the robustness check below\\
\bottomrule
\end{tabularx}}

\subsubsection{c) Result}
\res{$+22.3$ runs (190.4 vs 168.1)}{$+16.9$ to $+27.7$}{$t = 8.13$ (544)}{$1.5 \times 10^{-15}$ / $7.6 \times 10^{-15}$}{Cohen's $d = 0.68$}{reject $H_0$}
The whole interval clears 10 runs, so the rise \textbf{matters}. India-only matches (no UAE seasons): $+20.1$ (14.3 to 25.8).

\textbf{Rule or trend?} Totals were already rising by 0.94 runs a season. Net of that trend the Impact era adds \textbf{+18.1 runs} (11.9 to 24.3; season-clustered SE: 10.6 to 25.6), and comparing the same venues (fixed effects) gives $+19.5$ (12.0 to 27.0). Every lower bound still clears 10 runs. Placebo rules dated 2011--2021 never produce a jump above 7.7 runs (Fig.~\ref{fig:h1}), and the pre-2023 trend forecasts 170--173 for 2023--26 against observed means of 183--196.
\fig{fig5_h1_trend.pdf}{\label{fig:h1}Left: the pre-2023 trend (line, dashed when projected) badly under-predicts 2023--26 (blue). Right: estimated jump for placebo start years (grey) vs the real 2023 rule (blue), with 95\% CIs; dotted line = 10-run threshold.}

\subsection{H2 (P1): do teams hit more sixes per innings?}
\subsubsection{a) Hypotheses, written before running the test}
\hyp{$\mu_I - \mu_P \le 0$}{$\mu_I - \mu_P > 0$}{$\alpha = 0.05$}{one-sided (right)}{Welch t-test}{1 six per innings}
$\mu_I, \mu_P$: mean sixes per first innings, Impact and Pre eras. One-sided because P1 claims a rise. One six (6 runs) is the smallest change a viewer or analyst would notice.

\subsubsection{b) Other tests we considered}
Two independent groups and an unknown $\sigma$ call for a two-sample t-test. Six counts are discrete and mildly right-skewed (skew 0.41 and 0.48), but with 279 and 303 innings the sample means are close to normal by the central limit theorem. We use Welch's version because the spreads differ (SDs 4.34 and 3.72; Levene $p = 0.008$).
{\small
\begin{tabularx}{\columnwidth}{@{}P{2.3cm}L@{}}
\toprule
\textbf{Also considered} & \textbf{Why we did or did not use it}\\
\midrule
Pooled t-test & assumes equal variances, rejected above; $t = 6.85$, same decision\\
Mann--Whitney U & rank-based, no normality needed; run as the check: median 9 vs 7, HL shift $+2$ (1 to 3), $p = 2 \times 10^{-10}$, same decision\\
Poisson rate test & invalid: the variance is $\sim 2\times$ the mean (overdispersion), so its $z = 9.65$ is far too optimistic\\
Negative binomial regression & handles overdispersion and innings length (offset); run as the robustness check\\
\bottomrule
\end{tabularx}}

\subsubsection{c) Result}
\res{$+2.29$ sixes (9.28 vs 6.99)}{$+1.63$ to $+2.95$}{$t = 6.80$ (550)}{$1.3 \times 10^{-11}$ / $5.3 \times 10^{-11}$}{Cohen's $d = 0.57$}{reject $H_0$}
The whole interval clears 1 six, so the rise \textbf{matters}. Per 120 legal balls the medians are still 9 vs 7. The negative-binomial rate ratio is \textbf{1.33} (1.22 to 1.44): a third more sixes per ball; net of the trend it is 1.18 (1.06 to 1.30, $p = 0.002$). Sixes now supply 29\% of first-innings runs, up from 25\%.
\fig{fig6_sixes.pdf}{Six counts per first innings: the Impact-era distribution has a heavier right tail; innings with 13+ sixes rose from 8.6\% to 22.2\%.}

\subsection{H3 (P2): does winning the toss help win the match?}
\subsubsection{a) Hypotheses, written before running the test}
\hyp{$p_{\text{toss}} \le 0.5$}{$p_{\text{toss}} > 0.5$}{$\alpha = 0.05$}{one-sided (right)}{one-proportion z-test}{5 percentage points (55\%)}
$p_{\text{toss}}$: share of matches won by the toss winner; 1,218 matches with a winner (ties and no-results left out).

\subsubsection{b) Other tests we considered}
{\small
\begin{tabularx}{\columnwidth}{@{}P{2.3cm}L@{}}
\toprule
\textbf{Also considered} & \textbf{Why we did or did not use it}\\
\midrule
Exact binomial test & exact at any sample size, but unnecessary when $np = 609$; $p = 0.14$, same decision\\
Chi-square goodness of fit & equivalent but two-sided only ($\chi^2 = 1.19$, 1 d.f.)\\
TOST equivalence test & answers ``is the effect smaller than 5 pp?'', which a non-rejection cannot; run below\\
\bottomrule
\end{tabularx}}

\subsubsection{c) Result}
\res{51.6\% (628 of 1,218)}{48.8\% to 54.4\%}{$z = 1.09$}{0.14 / 0.28}{$+1.6$ pp; Cohen's $h = 0.03$}{fail to reject $H_0$}
Not rejecting is not the same as ``no effect'', so we tested equivalence: both one-sided z-tests reject ($p = 0.008$ against $p \ge 0.55$; $p < 10^{-5}$ against $p \le 0.45$). The toss effect is \textbf{equivalent to zero within $\pm$5 pp}; power to detect a 5 pp edge was 97\%. The \emph{decision} matters instead: toss winners who fielded won 9.4 pp more often than those who batted (3.4 to 15.3; odds ratio 1.48, $p = 0.002$, adjusted for era).
\fig{fig7_toss.pdf}{Toss winners' win share by era and decision (95\% exact CIs). The overall estimate (diamond) sits inside the $\pm$5 pp band; fielding first is ahead of batting first in every era.}

\subsection{H4 (P2): has the chasing edge changed across eras?}
\subsubsection{a) Hypotheses, written before running the test}
\hyp{$p_{E} = p_{P} = p_{I}$}{at least one era differs}{$\alpha = 0.05$}{two-sided (any direction)}{chi-square test of independence ($3 \times 2$)}{10 points between any two eras}
$p$: share of clean matches won by the chasing side in 2008--17, 2018--22 and 2023--26. Two-sided, because higher targets could make chasing either harder or easier.

\subsubsection{b) Other tests we considered}
{\small
\begin{tabularx}{\columnwidth}{@{}P{2.3cm}L@{}}
\toprule
\textbf{Also considered} & \textbf{Why we did or did not use it}\\
\midrule
Pairwise two-proportion z-tests & three tests need a correction and do not test ``any difference''; all Holm $p = 1$\\
Fisher's exact & meant for sparse tables; the smallest expected count here is 127\\
G-test & likelihood-ratio version of the same test: $G = 0.44$, same answer\\
Logistic regression on target and era & asks ``at the same target?'', which the table cannot; run below\\
\bottomrule
\end{tabularx}}

\subsubsection{c) Result}
\res{chase win share 54.4\%, 55.8\%, 53.0\%}{Impact $-$ Pre: $-2.7$ points ($-10.8$ to $+5.4$)}{$\chi^2 = 0.44$ (2)}{0.80 / 0.80}{Cramér's $V = 0.02$}{fail to reject $H_0$}
Chasing keeps a steady edge of about 4 points over a coin flip. \textbf{Why has it not moved, when targets rose 22 runs?} Each extra 10 runs of target cuts the odds of a successful chase by 30\% (OR 0.70, 0.66 to 0.73). But at the \emph{same} target, Impact-era chasers have about twice the odds (OR 1.97, 1.34 to 2.89, $p < 0.001$): chasing 181 succeeds 63\% of the time now against 47\% before. Better chasing exactly offsets higher targets.
\fig{fig8_chase_target.pdf}{Chance of a successful chase against the target (logistic fit; dots are binned observed shares, size $\propto n$). The Impact-era curve sits higher at every target, but its mean target (dotted) has moved right.}

\subsection{H5 (P3): does a powerplay collapse sink the batting-first team?}
\subsubsection{a) Hypotheses, written before running the test}
\hyp{$p_{C} - p_{N} \ge 0$}{$p_{C} - p_{N} < 0$}{$\alpha = 0.05$}{one-sided (left)}{two-proportion z-test}{15 percentage points}
$p_C, p_N$: win share of batting-first sides that did / did not lose 3+ wickets in overs 1--6. One first innings per match, so observations are independent.

\subsubsection{b) Other tests we considered}
{\small
\begin{tabularx}{\columnwidth}{@{}P{2.3cm}L@{}}
\toprule
\textbf{Also considered} & \textbf{Why we did or did not use it}\\
\midrule
Chi-square ($2\times2$) & the same test squared, two-sided only ($\chi^2 = 39.8$)\\
Fisher's exact & unnecessary at $n = 1{,}183$; OR $= 0.32$, $p \approx 8 \times 10^{-11}$\\
Logistic regression on powerplay wickets & uses the full count, adjusts for venue and trend, and tests an era interaction; run below\\
\bottomrule
\end{tabularx}}

\subsubsection{c) Result}
\res{$-25.8$ points (23.6\% vs 49.4\%)}{$-32.8$ to $-18.8$}{$z = -6.31$}{$1.4 \times 10^{-10}$ / $4.2 \times 10^{-10}$}{odds ratio 0.32}{reject $H_0$}
\figw{0.86}{fig9_powerplay_era.pdf}{Win share by powerplay wickets before and during the Impact era (95\% CIs). The slopes are parallel: the rule has not changed what an early collapse costs.}
A collapse halves the batting-first side's chances; the whole interval lies beyond $-15$ points, so it \textbf{matters}. Adjusted for venue and trend, each powerplay wicket multiplies the odds of winning by 0.61 (0.54 to 0.69), through a lower total (187 runs with no early wicket, 151 with three). Deeper batting does \textbf{not} soften it: the wickets $\times$ Impact interaction is OR 1.07 (0.82 to 1.39, $p = 0.62$). Chasing sides pay the same price (25.1\% vs 61.1\%).

\end{multicols}
\vspace{-2pt}
\noindent\begin{minipage}{\textwidth}
{\large\bfseries All five tests at a glance}\par\vspace{3pt}
{\small
\begin{tabularx}{\textwidth}{@{}P{1.9cm}P{3.6cm}P{1.6cm}LP{2.1cm}@{}}
\toprule
\textbf{Test} & \textbf{Main estimate (95\% CI)} & \textbf{p (Holm)} & \textbf{Strongest robustness check} & \textbf{Decision}\\
\midrule
H1 totals & $+22.3$ runs (16.9, 27.7) & $8 \times 10^{-15}$ & $+18.1$ runs beyond the pre-existing trend (11.9, 24.3); same venues $+19.5$; placebo years never exceed 7.7 & reject; matters\\
H2 sixes & $+2.29$ sixes (1.63, 2.95) & $5 \times 10^{-11}$ & per ball and net of trend: rate ratio 1.18 (1.06, 1.30) & reject; matters\\
H3 toss & 51.6\% (48.8, 54.4) & 0.28 & equivalence test: within $\pm$5 pp of 50\% ($p = 0.008$); the bat/field choice matters instead & fail to reject; no effect\\
H4 chase & 54.4 / 55.8 / 53.0\% & 0.80 & at the same target, Impact-era chasers' odds $\times$1.97 (1.34, 2.89): better chasing offsets higher targets & fail to reject; explained\\
H5 collapse & $-25.8$ points ($-32.8$, $-18.8$) & $4 \times 10^{-10}$ & venue + trend adjusted OR 0.61 per wicket (0.54, 0.69); no change in the Impact era ($p = 0.62$) & reject; matters\\
\bottomrule
\end{tabularx}}
\captionof{table}{\label{tab:sum}Every decision survives its strongest alternative explanation and the Holm correction. The three rejections clear their real-world thresholds; the two non-rejections are backed by an equivalence test (H3) and a mechanism (H4).}
\end{minipage}
\vspace{2pt}
\begin{multicols}{2}
\raggedcolumns
\section{Limitations}
\begin{itemize}
\item \textbf{Era, not only rule.} The Impact era also brought the 2024 two-bouncer rule, new squads and bats. Trend adjustment, venue fixed effects and placebo years narrow this down but cannot rule out a change that happened in 2023 for another reason.
\item \textbf{Not independent within a season.} Matches in one season share conditions; season-clustered errors widen the H1 interval (10.6 to 25.6) but do not change the decision.
\item \textbf{Venue mix.} 2009 was in South Africa, 2014 began and 2020--21 were partly in the UAE; India-only and venue fixed-effect checks keep H1.
\item \textbf{Matches left out.} 60 shortened, tied or abandoned matches (4.8\%) are not in the score and chase tests.
\item \textbf{Self-selection.} Captains choose to bat or field (H3), and collapses are more likely for weaker sides or on harder pitches (H5); these are associations.
\item \textbf{One league.} Results describe the IPL, not T20 cricket in general; four Impact seasons is still a short run.
\end{itemize}

\end{multicols}
\end{document}
